\chapter*{Introduction Générale}
\addcontentsline{toc}{chapter}{Introduction Générale}

Quand la pression monte dans le monde du travail, suivre sa clientèle devient central. Surtout en B2B. Ce n’est plus seulement un carnet d’adresses électronique. Ces systèmes touchent maintenant chaque moment de la vente. L’acheteur s’y sent mieux accompagné. Les chiffres finaux montrent souvent une amélioration. Pourtant, en regardant bien les options offertes, on tombe sur un même problème. Les petites équipes sont régulièrement mises de côté. Trop complexes parfois. Parfois mal conçus. Le fonctionnement paraît brouillon. Le prix fait peur. La vitesse chute avec beaucoup de données. L’ajustement aux habitudes locales est rarement fluide.

\vspace{0.5cm}

Ça commence par une idée simple. Un truc monté à partir de presque rien. Ce projet s'inscrit ici, non pas en cassant tout, mais en suivant le courant. Le but ? Quelques mots suffisent. Créer un CRM micro-SaaS solide, fait pour les petites équipes commerciales ciblant les pros. Depuis un moment déjà, on voit ce trou. Rien d’offert qui soit léger, vif, sans frais cachés, tout en restant puissant. Ce projet vise quelque chose de précis. Il veut mélanger puissance calme des grands outils et agilité des applications modernes. Ici, une interrogation monte. Comment créer un CRM fiable, qui ne bloque pas quand l’équipe grossit ? Un système ouvert aux petites équipes commerciales. Quelque chose de rapide même sous charge. Sans sacrifier le confort d’utilisation chaque jour.

\vspace{0.5cm}

C’est avant tout un CRM léger fait pour les petites équipes en B2B. Son cœur fonctionne fort, prêt à monter en puissance si besoin. Chaque détail tient la route, surtout le suivi des contacts qui glisse sans accroc. Pas de fouillis partout, l’écran montre juste l’essentiel. Derrière ce calme visuel se cache une machine bien réglée. Des transferts fluides circulent d’un système à l’autre grâce à des outils d’import modulables. En temps réel, chaque modification apparaît dès qu’elle est enregistrée. La sécurité veille en silence durant toutes les phases actives. C’est là un projet imposé dans le cadre universitaire pour clore le parcours ingénieur logiciel. Ce travail a pris forme sur une demi-année au sein de VisioAd, accompli en fin de cursus.

\vspace{0.5cm}

Tout débute comme ça. Un aperçu posé dès l'entrée pour donner une direction claire à ce qui suit. Le texte progresse étape après étape, sans se précipiter. Chaque partie examine un angle nouveau - d’abord le cadre du projet, ensuite les besoins repérés, puis une feuille de route réaliste, la phase d’exécution, enfin une évaluation globale avec quelques perspectives envisageables. Le départ se fait depuis ce qui existe déjà. Cela permet de s’appuyer sur des faits tangibles. Petit à petit, on y voit plus net dans les ressorts techniques, humains et fonctionnels liés à notre CRM Micro-SaaS.